% from ega0-17-fr.tex
\oldpage[0\textsubscript{IV-1}]{41}
En effet, l'idéal maximal de $\widehat{A}$ est $\mathfrak{m}\widehat{A}$ et on
sait que $\mathfrak{m}^h/\mathfrak{m}^{h+1}$ et
$(\mathfrak{m}\widehat{A})^h/(\mathfrak{m}\widehat{A})^{h+1}$ sont des
$k$-espaces vectoriels isomorphes~; la condition a) de (17.1.1) est donc la
même pour $A$ et $\widehat{A}$.
\end{proof}

\begin{definition}[17.1.6]
\label{0.17.1.6-fr}
Étant donné un anneau local noethérien $A$, d'idéal maximal $\mathfrak{m}$, on
appelle \emph{système régulier de paramètres de $A$} un système de paramètres
pour $A$ (16.3.5) qui engendre $\mathfrak{m}$.
\end{definition}

L'existence d'un tel système équivaut au fait que $A$ est régulier (17.1.1)~;
si $(x_i)_{1\leq i\leq n}$ est un système régulier de paramètres, c'est un
système minimal de générateurs de $\mathfrak{m}$, dont les classes mod.
$\mathfrak{m}^2$ forment une base de $\mathfrak{m}/\mathfrak{m}^2$ sur $k$~;
en vertu de (17.1.1, a)) et de (15.1.9), $(x_i)$ est une suite $A$-régulière
maximale (d'où la terminologie).

On prendra garde toutefois que, dans un anneau régulier, un système de
paramètres pour $A$ qui est aussi une suite $A$-régulière \emph{n'est pas
nécessairement} un système régulier de paramètres, comme le montre l'exemple
des puissances $k$-ièmes d'un système régulier de paramètres (15.1.20).

\begin{proposition}[17.1.7]
\label{0.17.1.7-fr}
Soient $A$ un anneau local noethérien, $\mathfrak{m}$ son idéal maximal,
$k=A/\mathfrak{m}$ son corps résiduel, $(x_i)_{1\leq i\leq r}$ une suite
d'éléments de $\mathfrak{m}$, $\mathfrak{J}$ l'idéal
$x_1A+\cdots+x_rA$. Les conditions suivantes sont équivalentes~:
\begin{enumerate}
  \item[\rm a)] $A$ est régulier et les $x_i$ font partie d'un système
  régulier de paramètres de $A$.
  \item[\rm a')] $A$ est régulier et les classes $\overline{x}_i$ des $x_i$
  mod. $\mathfrak{m}^2$ sont linéairement indépendantes sur $k$.
  \item[\rm b)] Les $x_i$ font partie d'un système de paramètres pour $A$ et
  $A/\mathfrak{J}$ est un anneau régulier.
\end{enumerate}
En outre, lorsque ces conditions sont vérifiées, $\mathfrak{J}$ est un idéal
premier.
\end{proposition}

\begin{proof}
La dernière assertion résulte trivialement du fait que $A/\mathfrak{J}$,
étant régulier, est intègre (17.1.3). Posons $n=\dim(A)$. L'équivalence de a)
et a') est immédiate~; en effet, les classes mod. $\mathfrak{m}^2$ des éléments
d'un système régulier de paramètres forment une base de
$\mathfrak{m}/\mathfrak{m}^2$ par rapport à $k$, et réciproquement, si les
$\overline{x}_i$ $(1\leq i\leq r)$ sont linéairement indépendants sur $k$, on
peut trouver $n-r$ éléments $x_i\in\mathfrak{m}$ $(r+1\leq i\leq n)$ tels que
les $\overline{x}_i$ $(1\leq i\leq n)$ forment une base de
$\mathfrak{m}/\mathfrak{m}^2$, donc $(x_i)_{1\leq i\leq n}$ est un système
régulier de paramètres. Pour voir que a) entraîne b), remarquons que puisque
les $x_i$ font partie d'un système de paramètres de $A$, on a
$\dim(A/\mathfrak{J})=n-r$ (16.3.6)~; soit
$\mathfrak{n}=\mathfrak{m}/\mathfrak{J}$ l'idéal maximal de
$A/\mathfrak{J}$. On a une suite exacte
\begin{equation}
\label{0.17.1.7.1-fr}
\tag{17.1.7.1}
  0\longrightarrow(\mathfrak{m}^2+\mathfrak{J})/\mathfrak{m}^2
  \longrightarrow\mathfrak{m}/\mathfrak{m}^2
  \longrightarrow\mathfrak{n}/\mathfrak{n}^2\longrightarrow0
\end{equation}
puisque $\mathfrak{n}^2=(\mathfrak{m}^2+\mathfrak{J})/\mathfrak{J}$~;
l'hypothèse a) entraîne que
$\operatorname{rg}_k((\mathfrak{m}^2+\mathfrak{J})/\mathfrak{m}^2)=r$, donc
$\operatorname{rg}_k(\mathfrak{n}/\mathfrak{n}^2)=n-r
=\dim(A/\mathfrak{J})$, ce qui prouve b) en vertu de (17.1.1, b)).

Inversement, pour prouver que b) entraîne a), notons que le fait que les $x_i$
font partie d'un système de paramètres entraîne que
$\dim(A/\mathfrak{J})=n-r$ (16.3.6)~; l'hypothèse que $A/\mathfrak{J}$ est
régulier implique d'autre part
$\operatorname{rg}_k(\mathfrak{n}/\mathfrak{n}^2)=n-r$ (17.1.1)~; par ailleurs
on a évidemment
$\operatorname{rg}_k((\mathfrak{m}^2+\mathfrak{J})/\mathfrak{m}^2)\leq r$,
donc on tire de (17.1.7.1) que
$\operatorname{rg}_k(\mathfrak{m}/\mathfrak{m}^2)\leq r+(n-r)=n$~; donc $A$
est régulier en vertu de (16.2.6) et (17.1.1). En outre la relation
$\operatorname{rg}_k(\mathfrak{m}/\mathfrak{m}^2)=n$ entraîne alors
$\operatorname{rg}_k((\mathfrak{m}^2+\mathfrak{J})/\mathfrak{m}^2)=r$, et par
suite les $\overline{x}_i$ sont linéairement indépendants sur $k$. C.Q.F.D.
\end{proof}

